\chapter[PVD in Framework Analysis]{Principle of Virtual Displacements in Framework Analysis}
\label{ch:pvd2}

\section{Shape functions and stiffness matrix}
\subsection{Two-noded (linear) axial element}
\begin{equation}
  u(x) = \Bigl(1-\frac xL\Bigr)u_1 + \frac xL u_2 = N_1(x)u_1 + N_2(x)u_2 = \langle N_1, N_2\rangle\Bmat{u_1\\ u_2} = \NN\bs{\Delta}
\end{equation}
where the shape functions are
\begin{equation}
  N_1(x) = 1-\frac xL\quad\text{and}\quad N_2(x) = \frac xL
\end{equation}
Note that shape function $N_i(x)$ represents the displacement field associated with $\Delta_i = 1$ and other $\Delta_i$'s $= 0$.

The strain field is then given by
\begin{equation}
  \eps(x) = u'(x) = \NN'\bs{\Delta}
\end{equation}
where $\NN' = \langle N_1', N_2'\rangle = \langle -1/L, 1/L\rangle$.

The virtual strain energy is
\begin{equation}
  \dU = \int_0^L\delta\eps\,E\eps\,A\,dx = \int_0^L(\delta\eps)^\tr E\eps\,A\,dx = \delta\bs{\Delta}^\tr\Bigl(\int_0^L EA\NN'^{\tr}\NN'\,dx\Bigr)\bs{\Delta}
\end{equation}
Noting that $\dU = \delta\bs{\Delta}^\tr\kmat\bs{\Delta}$, therefore,
\begin{gather}
  \kmat = \int_0^L EA\NN'^{\tr}\NN'\,dx = EA\int_0^L\NN'^{\tr}\NN'\,dx\quad\text{if } EA = \text{constant}\\
  \Rightarrow\quad \kmat = \frac{EA}{L}\bmat{1 & -1\\ -1 & 1}\quad\text{if } EA = \text{constant.}
\end{gather}

\subsection{Flexural element: bent about the \texorpdfstring{$z$}{z}-axis}
The displacement is
\begin{equation}
  v(x) = N_1(x)v_1 + N_2(x)\theta_{z1} + N_3(x)v_2 + N_4(x)\theta_{z2} = \sum_{i=1}^4 N_i(x)\Delta_i = \NN\bs{\Delta}
\end{equation}
where the shape functions are
\begin{align*}
  N_1(x) &= 2\Bigl(\frac xL\Bigr)^3 - 3\Bigl(\frac xL\Bigr)^2 + 1\\
  N_2(x) &= \frac{x^3}{L^2} - 2\frac{x^2}{L} + x\\
  N_3(x) &= -2\Bigl(\frac xL\Bigr)^3 + 3\Bigl(\frac xL\Bigr)^2\\
  N_4(x) &= \frac{x^3}{L^2} - \frac{x^2}{L}
\end{align*}
Similarly, shape function $N_i(x)$ is the displacement field due to $\Delta_i = 1$ and the other $\Delta_i$'s $= 0$. For example, $N_2(x)$ satisfies the following four conditions: $N_2(0) = 0$, $N_2'(0) = 1$, $N_2(L) = 0$, $N_2'(L) = 0$; these correspond to a hinged support with a unit rotation for the left end and a fixed support for the right end.

The above shape functions can be derived as follows. First, as there are four degrees of freedom, a cubic displacement field can be assumed; namely
\begin{equation}
  v(x) = a_0 + a_1x + a_2x^2 + a_3x^3
\end{equation}
Then use of the four boundary conditions: $v(0) = v_1$, $v(L) = v_2$, $v'(0) = \theta_{z1}$, and $v'(L) = \theta_{z2}$ results in
\begin{equation}
  \bmat{1 & 0 & 0 & 0\\ 1 & L & L^2 & L^3\\ 0 & 1 & 0 & 0\\ 0 & 1 & 2L & 3L^2}\Bmat{a_0\\ a_1\\ a_2\\ a_3} = \Bmat{v_1\\ v_2\\ \theta_{z1}\\ \theta_{z2}}
  \quad\Rightarrow\quad
  \begin{aligned}
    a_0 &= v_1\\ a_1 &= \theta_{z1}\\ a_2 &= \frac{3}{L^2}(-v_1+v_2) - \frac1L(2\theta_{z1}+\theta_{z2})\\ a_3 &= \frac{2}{L^3}(v_1-v_2) + \frac{1}{L^2}(\theta_{z1}+\theta_{z2})
  \end{aligned}
\end{equation}
Substituting $a_0, a_1, a_2$, and $a_3$ into $v(x)$ and rearranging the terms results in the shape functions immediately.

The axial strain along the depth of a cross section due to bending is given by
\begin{equation}
  \eps(y) = -yv'' = -y\NN''\bs{\Delta}
\end{equation}
The virtual strain energy is
\begin{equation}
  \begin{aligned}
    \dU &= \int_V\delta\eps\,\sigma\,dV = \int_V\delta\eps\,E\eps\,dV\\
        &= \int_V\delta v''Ey^2v''\,dV = \int_0^L EI_z\,\delta v''^{\tr}v''\,dx\\
        &= \delta\bs{\Delta}^\tr\Bigl(\int_0^L EI_z\NN''^{\tr}\NN''\,dx\Bigr)\bs{\Delta} = \delta\bs{\Delta}^\tr\kmat\bs{\Delta}
  \end{aligned}
\end{equation}
Therefore,
\begin{equation}
  \kmat = \int_0^L EI_z\NN''^{\tr}\NN''\,dx\qquad\text{or}\qquad k_{ij} = \int_0^L EI_zN_i''(x)N_j''(x)\,dx
\end{equation}
If $EI_z =$ constant, it is easy to obtain, on the basis of this equation and the shape functions, that
\begin{equation}
  \kmat = \frac{EI_z}{L}
  \begin{bmatrix}
    \dfrac{12}{L^2} & \dfrac6L & -\dfrac{12}{L^2} & \dfrac6L\\[8pt]
    & 4 & -\dfrac6L & 2\\[8pt]
    & & \dfrac{12}{L^2} & -\dfrac6L\\[8pt]
    \multicolumn{2}{l}{\textit{Symm.}} & & 4
  \end{bmatrix}
\end{equation}

\subsection{Flexural element: bent about the \texorpdfstring{$y$}{y}-axis}
The displacement field is:
\begin{equation}
  w(x) = N_1(x)w_1 + N_2(x)\theta_{y1} + N_3(x)w_2 + N_4(x)\theta_{y2} = \sum_{i=1}^4 N_i(x)\Delta_i = \NN\bs{\Delta}
\end{equation}
where the shape functions are given by
\begin{align*}
  N_1(x) &= 2\Bigl(\frac xL\Bigr)^3 - 3\Bigl(\frac xL\Bigr)^2 + 1\\
  N_2(x) &= -\Bigl(\frac{x^3}{L^2} - 2\frac{x^2}{L} + x\Bigr)\\
  N_3(x) &= -2\Bigl(\frac xL\Bigr)^3 + 3\Bigl(\frac xL\Bigr)^2\\
  N_4(x) &= -\Bigl(\frac{x^3}{L^2} - \frac{x^2}{L}\Bigr)
\end{align*}
Note that shape functions $N_1(x)$ and $N_3(x)$ are the same as those for the $z$-axis bending but shape functions $N_2(x)$ and $N_4(x)$ have a sign change compared with those for the $z$-axis bending. The sign change is due to the fact that $\theta_y(x) = -w'(x)$ but $\theta_z(x) = v'(x)$.

\subsection{Torsional element}
The shape functions for $\theta_x$ are the same as those for the axial displacement. Therefore, the stiffness matrix for a torsional element takes the form:
\begin{equation}
  \kmat = \int_0^L GJ\NN'^{\tr}\NN'\,dx \quad\Rightarrow\quad \kmat = \frac{GJ}{L}\bmat{1 & -1\\ -1 & 1}\ \text{if } GJ = \text{constant.}
\end{equation}

\section{Special purpose elements}
\subsection{Three-noded (quadratic) axial element}
The displacement field is given by
\begin{equation}
  u(x) = \sum_{i=1}^3 N_i(x)u_i = \NN(x)\bs{\Delta}
\end{equation}
where $\NN(x) = \langle N_1(x), N_2(x), N_3(x)\rangle$ and $\bs{\Delta} = \langle u_1, u_2, u_3\rangle^\tr$. The shape functions are given by
\begin{fig}
  % Shape functions of the three-noded (quadratic) axial element
\begin{tikzpicture}[scale=1.0,font=\small,
  declare function={Na(\x)=1-3*\x+2*\x*\x; Nb(\x)=4*\x-4*\x*\x; Nc(\x)=-\x+2*\x*\x;}]
  \def\L{6}
  % N1
  \begin{scope}
    \draw (0,0) -- (\L,0);
    \draw[member,domain=0:1,samples=60,variable=\t] plot ({\t*\L},{Na(\t)*1.2});
    \draw[dim] (-0.3,0) -- (-0.3,1.2) node[midway,left] {1};
    \node[below] at (0,-0.05) {0}; \node[below] at (\L/2,-0.05) {$\frac L2$}; \node[below] at (\L,-0.05) {$L$};
    \node[right] at (\L+0.4,0.6) {$N_1(x) = \dfrac{2}{L^2}\Bigl(x-\dfrac L2\Bigr)(x-L) = 1-3\dfrac xL+2\dfrac{x^2}{L^2}$};
  \end{scope}
  % N2
  \begin{scope}[yshift=-2.6cm]
    \draw (0,0) -- (\L,0);
    \draw[member,domain=0:1,samples=60,variable=\t] plot ({\t*\L},{Nb(\t)*1.2});
    \draw[dim] (\L/2,0) -- (\L/2,1.2) node[midway,right] {1};
    \node[below] at (0,-0.05) {0}; \node[below] at (\L/2,-0.05) {$\frac L2$}; \node[below] at (\L,-0.05) {$L$};
    \node[right] at (\L+0.4,0.6) {$N_2(x) = \dfrac{-4}{L^2}x(x-L) = 4\dfrac xL-4\dfrac{x^2}{L^2}$};
  \end{scope}
  % N3
  \begin{scope}[yshift=-5.2cm]
    \draw (0,0) -- (\L,0);
    \draw[member,domain=0:1,samples=60,variable=\t] plot ({\t*\L},{Nc(\t)*1.2});
    \draw[dim] (\L+0.3,0) -- (\L+0.3,1.2) node[midway,right] {1};
    \node[below] at (0,-0.05) {0}; \node[below] at (\L/2,-0.05) {$\frac L2$}; \node[below] at (\L,-0.05) {$L$};
    \node[right] at (\L+0.8,0.6) {$N_3(x) = \dfrac{2}{L^2}x\Bigl(x-\dfrac L2\Bigr) = -\dfrac xL+2\dfrac{x^2}{L^2}$};
  \end{scope}
\end{tikzpicture}

\end{fig}
Let the cross-sectional area be of the form $A(x) = A_0(1-mx/L)$, where $m =$ constant. The stiffness matrix for this linearly tapered axial bar is:
\begin{fig}
  % Linearly tapered axial bar with three nodes
\begin{tikzpicture}[scale=1.0,font=\small]
  \draw[thick] (0,-1.0) -- (6,-0.35) -- (6,0.35) -- (0,1.0) -- cycle;
  \draw[dashed] (-0.8,0) -- (7.2,0);
  \draw (3,-0.68) -- (3,0.68);
  \draw[load] (0,0) -- (0.9,0) node[below] {$u_1$};
  \draw[load] (3,0) -- (3.9,0) node[below] {$u_2$};
  \draw[load] (6,0) -- (6.9,0) node[below] {$u_3$};
  \node[above left] at (0,1.0) {$A_0$};
  \node[above] at (3,0.7) {$A_0\bigl(1-m\frac12\bigr)$};
  \node[above right] at (6,0.35) {$A_0(1-m)$};
  \node[below] at (0,-1.1) {$x=0$}; \node[below] at (3,-0.8) {$x=0.5L$}; \node[below] at (6,-0.45) {$x=L$};
\end{tikzpicture}

\end{fig}
\begin{equation}
  \kmat = \int_0^L EA\NN'^{\tr}\NN'\,dx = \frac{EA_0}{6L}\bmat{-3m+14 & 4m-16 & -m+2\\ 4m-16 & -16m+32 & 12m-16\\ -m+2 & 12m-16 & -11m+14}
\end{equation}
It is easy to verify that each column of $\kmat$ is a set of forces in equilibrium (summation of each column is zero). A special case is that of $A(x) = A_0$, i.e.\ $m = 0$, for which
\begin{equation}
  \kmat = \frac{EA_0}{3L}\bmat{7 & -8 & 1\\ -8 & 16 & -8\\ 1 & -8 & 7}
\end{equation}

\subsection{Exact shape functions and exact stiffness matrix for a tapered bar}
The governing differential equation for a tapered bar is given by
\begin{equation}
  \od{}{x}\Bigl(EA(x)\od{u}{x}\Bigr) = 0
\end{equation}
where $A(x) = A_0(1-mx/L)$ is assumed. Integrating the differential equation twice and using the boundary conditions $u(0) = u_1$ and $u(L) = u_2$ yields
\begin{align*}
  &\od{}{x}\Bigl(EA_0(1-mx/L)\od{u}{x}\Bigr) = 0\\
  &\Rightarrow (1-mx/L)\od{u}{x} = c_1 \quad\Rightarrow\quad \od{u}{x} = \frac{c_1}{(1-mx/L)}\\
  &\Rightarrow u(x) = c_1\Bigl(\frac{-L}{m}\Bigr)\ln(1-mx/L) + c_2\\
  &u(0) = u_1 \quad\Rightarrow\quad c_2 = u_1\\
  &u(L) = u_2 \quad\Rightarrow\quad c_1\Bigl(\frac{-L}{m}\Bigr)\ln(1-m) + u_1 = u_2\\
  &\Rightarrow c_1 = (u_2-u_1)\Bigl(\frac{-m}{L}\Bigr)\frac{1}{\ln(1-m)}\\
  &\Rightarrow u(x) = (u_2-u_1)\frac{\ln(1-mx/L)}{\ln(1-m)} + u_1 = N_1(x)u_1 + N_2(x)u_2
\end{align*}
where
\begin{equation}
  N_1(x) = 1 - \frac{\ln(1-mx/L)}{\ln(1-m)},\qquad N_2(x) = \frac{\ln(1-mx/L)}{\ln(1-m)}
\end{equation}
The above shape functions are exact in the sense that they satisfy the differential equation. The ``exact'' stiffness matrix, calculated by $k_{ij} = E\int_0^L A(x)N_i'(x)N_j'(x)\,dx$, is of the form:
\begin{equation}
  \kmat = -\frac{EA_0}{L}\frac{m}{\ln(1-m)}\bmat{1 & -1\\ -1 & 1}
\end{equation}
If a two-noded linear element is used, then $u(x) = (1-\frac xL)u_1 + \frac xL u_2 \Rightarrow \od{u}{x} = \dfrac{u_2-u_1}{L}$
\begin{equation*}
  \od{}{x}\Bigl[EA_0(1-mx/L)\Bigl(\frac{u_2-u_1}{L}\Bigr)\Bigr] \neq 0
\end{equation*}

\begin{example}
Consider a linearly tapered bar fixed at the left end and a force $P$ is applied at the right end. Assuming $m = 2/3$, the exact stiffness equations are:
\begin{equation*}
  -\frac{EA_0}{L}\frac{2/3}{\ln(1/3)}\bmat{1 & -1\\ -1 & 1}\Bmat{u_1\\ u_2} = \Bmat{R\\ P}
\end{equation*}
where $R$ is the reaction at the fixed end. Using the boundary condition $u_1 = 0$, we obtain from the above stiffness equations that
\begin{equation*}
  u_2 = 1.6479\frac{PL}{EA_0}
\end{equation*}
which is, of course, exact.

If a quadratic axial element is used to solve the problem, the stiffness equations can be written as
\begin{equation*}
  \frac{2EA_0}{9L}\bmat{9 & -10 & 1\\ -10 & 16 & -6\\ 1 & -6 & 5}\Bmat{u_1\\ u_2\\ u_3} = \Bmat{R\\ 0\\ P}
\end{equation*}
Knowing that $u_1 = 0$, we can solve for $u_2$ and $u_3$ from the above equations that
\begin{equation*}
  \Bmat{u_2\\ u_3} = \Bmat{0.6136\\ 1.6364}\frac{PL}{EA_0}
\end{equation*}
The error for $u_3$, compared with the exact value, is only 0.7\%.

If a two-noded linear element is used and the cross-sectional area is taken as the mean cross-sectional area, i.e.\ $\overline{A} = \frac23A_0$, then we have
\begin{equation*}
  \frac{2EA_0}{3L}\bmat{1 & -1\\ -1 & 1}\Bmat{u_1\\ u_2} = \Bmat{R\\ P} \Rightarrow u_2 = 1.5\frac{PL}{EA_0}
\end{equation*}
The error is greater than that with the quadratic element. If the bar is divided into three or more elements, the error will become very small.
\end{example}

\subsection{Exact stiffness matrix of a tapered beam}
\begin{minipage}{0.45\linewidth}
Assume a tapered beam with width $b =$ constant and depth $h(x) = h_1\bigl(1+\frac xL\bigr)$. Therefore
\begin{equation*}
  I(x) = \frac{bh^3}{12} = \frac{bh_1^3}{12}\Bigl(1+\frac xL\Bigr)^3
\end{equation*}
\end{minipage}\hfill
\begin{minipage}{0.5\linewidth}\centering
  % Tapered beam: depth h(x) = h1(1 + x/L)
\begin{tikzpicture}[scale=1.0,font=\small]
  \draw[thick] (0,-0.5) -- (6,-1.0) -- (6,1.0) -- (0,0.5) -- cycle;
  \draw[dim] (3,-0.75) -- (3,0.75); \node[right] at (3.1,0.3) {$h(x)$};
  \node[left] at (-0.1,0) {$h_1$}; \node[right] at (6.1,0) {$2h_1$};
  \node[below] at (0,-0.6) {$x=0$}; \node[below] at (6,-1.1) {$x=L$};
\end{tikzpicture}

\end{minipage}

Governing equation (G.E.):
\begin{equation}
  \frac{d^2}{dx^2}\bigl[EI(x)v''\bigr] = 0
\end{equation}
By integrating G.E. twice, we have
\begin{equation}
  \Bigl(1+\frac xL\Bigr)^3v'' = c_1 + c_2x
\end{equation}
By integrating the above equation twice, we obtain
\begin{equation}
  v(x) = \frac{L^2}{2}\Bigl(\frac{c_1}{1+\frac xL}\Bigr) - L^3\Bigl[\frac{1}{2(1+\frac xL)} + \ln\Bigl(1+\frac xL\Bigr)\Bigr]c_2 + c_3x + c_4
\end{equation}
Boundary conditions (BCs):
\begin{equation*}
  v(0) = v_1,\quad v'(0) = \theta_{z1},\quad v(L) = v_2,\quad v'(L) = \theta_{z2}
\end{equation*}
Substituting BCs into $v(x)$, we can obtain $c_i$ ($i = 1, 4$) in terms of $v_1, \theta_{z1}, v_2$, and $\theta_{z2}$. Finally, rearrange the terms to result in the form below:
\begin{equation}
  v(x) = N_1(x)v_1 + N_2(x)\theta_{z1} + N_3(x)v_2 + N_4(x)\theta_{z2}
\end{equation}
The shape functions $N_i(x)$ ($i = 1, 4$) are therefore obtained; these shape functions are exact since they satisfy the governing equation.

Stiffness matrix:
\begin{equation}
  \kmat = \int_0^L EI(x)\NN''^{\tr}\NN''\,dx = \text{exact stiffness matrix.}
\end{equation}

\section{Use of PVD to calculate the equivalent nodal forces or fixed-end forces}
In general, we can determine $\mat{f}_{eq}$ using the principle of virtual displacements. In this way, we let the virtual work done by $\mat{f}_{eq}$ be equal to that done by the intermediate concentrated loads (can be expressed as delta functions) or distributed loads. That is,
\begin{fig}
  % Distributed load q(x) on a beam -> equivalent nodal loads
\begin{tikzpicture}[scale=0.9,font=\small]
  \begin{scope}
    \draw[thick] (0,0) rectangle (5,0.4);
    \draw[thick] (0,0.4) .. controls (0.8,2.2) and (1.6,2.0) .. (2.2,1.6) .. controls (2.8,1.2) and (3.6,1.4) .. (4.2,1.9) .. controls (4.6,2.2) .. (5,0.4);
    \foreach \x in {0.4,0.8,...,4.7}{\draw (\x,0.42) -- (\x,{0.4+1.3});}
    \node[above] at (2.6,1.9) {$q(x)$};
  \end{scope}
  \node at (6.2,0.6) {$\Longrightarrow$};
  \begin{scope}[xshift=7.5cm]
    \draw[thick] (0,0) rectangle (5,0.4);
    \draw[load] (0,0.5) -- (0,1.6); \draw[load] (5,0.5) -- (5,1.6);
    \draw[->,thick] (-0.3,-0.2) arc (-90:90:0.4); \draw[->,thick] (5.3,-0.2) arc (270:90:0.4);
  \end{scope}
\end{tikzpicture}

\end{fig}
\begin{gather}
  \delta\bs{\Delta}^\tr\mat{f}_{eq} = \int_0^L(\NN\delta\bs{\Delta})^\tr q(x)\,dx
  \;\rightarrow\; \delta\bs{\Delta}^\tr\mat{f}_{eq} = \delta\bs{\Delta}^\tr\int_0^L\NN^\tr q(x)\,dx\\
  \rightarrow\quad \mat{f}_{eq} = \int_0^L\NN^\tr q(x)\,dx
\end{gather}
where $\NN = \langle N_1, N_2, \dots\rangle$ are shape functions. In component form,
\begin{equation}
  (\mat{f}_{eq})_i = \int_0^L N_i(x)q(x)\,dx
\end{equation}
The corresponding fixed-end forces are $\mat{f}_{fef} = -\mat{f}_{eq}$.

\begin{example}
\leavevmode
\begin{fig}
  % Uniform load q -> equivalent nodal loads f_eq1..f_eq4
\begin{tikzpicture}[scale=0.9,font=\small]
  \begin{scope}
    \draw[thick] (0,0) rectangle (5,0.4);
    \foreach \x in {0.2,0.7,...,4.9}{\draw[->] (\x,0.45) -- (\x,1.3);}
    \node[above] at (2.5,1.3) {$q$};
  \end{scope}
  \node at (6.2,0.4) {$\Longrightarrow$};
  \begin{scope}[xshift=7.5cm]
    \draw[thick] (0,0) rectangle (5,0.4);
    \draw[load] (0,0.5) -- (0,1.6) node[above] {$f_{eq1}$}; \draw[load] (5,0.5) -- (5,1.6) node[above] {$f_{eq3}$};
    \draw[->,thick] (-0.3,-0.2) arc (-90:90:0.4); \node[below left] at (-0.1,-0.1) {$f_{eq2}$};
    \draw[->,thick] (5.3,-0.2) arc (270:90:0.4); \node[below right] at (5.1,-0.1) {$f_{eq4}$};
  \end{scope}
\end{tikzpicture}

\end{fig}
Using the equations given above, we obtain
\begin{alignat*}{2}
  f_{eq1} &= q\int_0^L N_1(x)\,dx = \frac{qL}{2};&\qquad f_{eq3} &= q\int_0^L N_3(x)\,dx = \frac{qL}{2}\\
  f_{eq2} &= q\int_0^L N_2(x)\,dx = \frac{qL^2}{12};&\qquad f_{eq4} &= q\int_0^L N_4(x)\,dx = -\frac{qL^2}{12}
\end{alignat*}
\end{example}

\begin{example}
\leavevmode
\begin{fig}
  % Point load P at x = a -> equivalent nodal loads
\begin{tikzpicture}[scale=0.9,font=\small]
  \begin{scope}
    \draw[thick] (0,0) rectangle (5,0.4);
    \draw[load] (1.5,0.45) -- (1.5,1.8) node[above] {$P$};
    \draw[dim] (0,-0.5) -- (1.5,-0.5) node[midway,below] {$a$};
    \draw[dim] (1.5,-0.5) -- (5,-0.5) node[midway,below] {$b$};
    \draw (0,-0.7) -- (0,-0.3); \draw (1.5,-0.7) -- (1.5,-0.3); \draw (5,-0.7) -- (5,-0.3);
  \end{scope}
  \node at (6.2,0.4) {$\Longrightarrow$};
  \begin{scope}[xshift=7.5cm]
    \draw[thick] (0,0) rectangle (5,0.4);
    \draw[load] (0,0.5) -- (0,1.6) node[above] {$f_{eq1}$}; \draw[load] (5,0.5) -- (5,1.6) node[above] {$f_{eq3}$};
    \draw[->,thick] (-0.3,-0.2) arc (-90:90:0.4); \node[below left] at (-0.1,-0.1) {$f_{eq2}$};
    \draw[->,thick] (5.3,-0.2) arc (270:90:0.4); \node[below right] at (5.1,-0.1) {$f_{eq4}$};
  \end{scope}
\end{tikzpicture}

\end{fig}
We can use the same equations to calculate the equivalent nodal loads. In this case, $q(x) = P\delta(x-a)$. By using the property $\int_0^L N_i(x)\delta(x-a)\,dx = N_i(a)$, one can obtain that $f_{eqi} = PN_i(a)$ and therefore
\begin{equation*}
  \mat{f}_{eq} = \Bmat{\dfrac{Pb^2}{L^3}(3a+b)\\[6pt] \dfrac{ab^2P}{L^2}\\[6pt] \dfrac{Pa^2}{L^3}(a+3b)\\[6pt] -\dfrac{ba^2P}{L^2}}
  \qquad\text{Check special case } a = b = \frac L2:\quad
  \mat{f}_{eq} = \Bmat{\dfrac P2\\[6pt] \dfrac{PL}{8}\\[6pt] \dfrac P2\\[6pt] -\dfrac{PL}{8}}
\end{equation*}
or
\begin{equation*}
  \dW = P\delta v(x)\big|_{x=a} = P\bigl[N_1(x)\big|_{x=a}\delta v_1 + \dots + N_4(x)\big|_{x=a}\delta\theta_2\bigr] = f_{eqi}\delta\Delta_i \Rightarrow f_{eq_i} = PN_i(a)
\end{equation*}
Similarly, for an applied moment, $M$, at $x = a$, we have
\begin{equation*}
  \dW = M\delta v'(x)\big|_{x=a} = M\bigl[N_1'(x)\big|_{x=a}\delta v_1 + \dots + N_4'(x)\big|_{x=a}\delta\theta_2\bigr] = f_{eqi}\delta\Delta_i \Rightarrow f_{eq_i} = MN_i'(a)
\end{equation*}
(see Fig.~\ref{fig:fef-table}).
\end{example}

\begin{example}[Axial element]
\leavevmode
\begin{fig}
  % Triangular axial load q(x) -> equivalent nodal loads
\begin{tikzpicture}[scale=0.9,font=\small]
  \begin{scope}
    \draw[thick] (0,0) rectangle (5,0.4);
    \draw[thick] (0,0.4) -- (5,1.9) -- (5,0.4);
    \node[above] at (5,1.9) {$q$}; \node[right] at (5.05,1.1) {$q(x)$};
    \foreach \x in {0.3,1.3,...,4.4}{\draw[->] (\x,0.2) -- (\x+0.7,0.2);}
  \end{scope}
  \node at (6.2,0.4) {$\Longrightarrow$};
  \begin{scope}[xshift=7.5cm]
    \draw[thick] (0,0) rectangle (5,0.4);
    \draw[load] (0.1,0.2) -- (0.9,0.2); \node[above] at (0.4,0.5) {$f_{eq1}$};
    \draw[load] (5.05,0.2) -- (5.9,0.2); \node[above] at (5.3,0.5) {$f_{eq2}$};
  \end{scope}
\end{tikzpicture}

\end{fig}
\begin{align*}
  f_{eq1} &= \int_0^L q(x)N_1(x)\,dx = \int_0^L q\frac xL\Bigl(1-\frac xL\Bigr)dx = \frac{qL}{6}\\
  f_{eq2} &= \int_0^L q(x)N_2(x)\,dx = \int_0^L q\frac xL\Bigl(\frac xL\Bigr)dx = \frac{qL}{3}
\end{align*}
Consider the axial bar below discretized into 2 elements. We can use the above formulas to obtain nodal loads as
\begin{equation*}
  \Pvec = \Pvec_c + \Pvec_{eq} = \Bmat{0\\ P} + \Bmat{\frac12qL\bigl(\frac23\bigr)+\frac12qL\bigl(\frac23\bigr)\\[4pt] \frac12qL\bigl(\frac13\bigr)} = \Bmat{\frac23qL\\[4pt] P+\frac16qL}
\end{equation*}
By $\Kmat\bs{\Delta} = \Pvec$, $\bs{\Delta}$ can be obtained.
\begin{fig}
  % Axial bar discretized into two elements with triangular load and a tip load P
\begin{tikzpicture}[scale=1.0,font=\small]
  \fill[left color=black!60,right color=black!15] (-0.3,-0.45) rectangle (0,0.45);
  \draw[thick] (0,-0.2) rectangle (6,0.2);
  \draw (3,-0.2) -- (3,0.2);
  \draw[thick] (0,0.2) -- (3,1.4) -- (6,0.2);
  \node[above] at (3,1.4) {$q$};
  \foreach \x in {0.3,1.2,...,5.4}{\draw[->] (\x,0) -- (\x+0.6,0);}
  \draw[load] (6.1,0) -- (7.0,0) node[right] {$P$};
  \node[above left] at (0.15,0.25) {1}; \node[above] at (3,0.25) {2}; \node[above right] at (5.85,0.25) {3};
  \draw[dim] (0,-0.8) -- (3,-0.8) node[midway,fill=white,inner sep=1pt] {$L$};
  \draw[dim] (3,-0.8) -- (6,-0.8) node[midway,fill=white,inner sep=1pt] {$L$};
  \foreach \x in {0,3,6}{\draw (\x,-0.6) -- (\x,-1.0);}
\end{tikzpicture}

\end{fig}
\end{example}
